\documentclass[11pt]{article}
\begin{document}
\section*{Step 41 Tensor-Network RT Bound}
This statement is scoped to the declared finite tensor-network carrier.

For each tested tensor network, the boundary state \(|\psi_{\partial}\rangle\) is obtained by explicitly contracting internal indices. For a boundary region \(A\), the reduced density matrix \(\rho_A\) is then formed from the contracted amplitude vector, and
\[
  S(A)=-\operatorname{Tr}(\rho_A\log\rho_A)
\]
is computed from its spectrum.

The geometric area is computed separately as a min-cut in the capacitated bulk graph:
\[
  \operatorname{Area}(\gamma_A)=\min_{\gamma:A|\bar A}\sum_{e\in\gamma} \log D_e.
\]
On all tested contracted states,
\[
  S(A)\le \operatorname{Area}(\gamma_A).
\]
The copy-tensor holographic cases saturate this bound at two bond dimensions. Product and weighted tensor controls produce strict gaps, so the equality is not an identity imposed by construction.

This finite saturable-bound structure recognition-lands on the Ryu--Takayanagi/random-tensor-network form. It is not a continuum coefficient derivation and not a quantum-gravity solution.
\end{document}
